\section{Project Types and Phases}


\begin{frame}{Projects}
    \mycomment{
    What qualifies as a project, and what doesn't? -- This may seem obvious, but it is not. Quite often simple tasks are already called projects. But even more common, a \textbf{goal} is confused with a project. A project is always a process: the compley process to achieve a specific goal.
    }
    \mytext{
    
    \begin{itemize}%[<+->]
        \item What is \emph{not }a project? $\rightarrow$ `business as usual'
        \item Projects have a beginning, a middle and an end...
        \item Someone \emph{wants }a project
        \item Project \emph{makes a difference} $\rightarrow$ change
        \item Projects have specifics
        \item Projects involve people and other resources
        \item Projects have sponsors, owners, managers...
        \item Projects can be small, large...
        \item Projects can be services, actions, products...
        \item Products can be tangible, intangible...
        \item Projects are planned...
    \end{itemize}
    }

\end{frame}

\begin{frame}{Example projects}
    \mytext{
    Think of one project example!
    }
\end{frame}
\begin{frame}{Example projects}
    \mycomment{
    Here are some examples for projects. You see that there are projects in many fields, sectors and industries. Many projects involve many different fields, persons, and industries. Project management appears in many disciplines, and integrates many professions.
    }
    \mytext{
    \begin{itemize}%[<+->]
        \item Office move
        \item Construction project
        \item Business expanse
        \item Business acquisition
        \item Conference, event
        \item IT project
        \item Moon landing
        \item Military projects
        \item Privatisation
        \item Founding of an institution
        \item Disbanding, dismantling
    \end{itemize}
    }
\end{frame}

\begin{frame}{Projects are \textbf{not}...}
    \mycomment{
    A vague idea or wish does not qualify as a project yet. For this, the actual planning process must already have proceeded to some degree. Requirements must have become clear:  specifics about the actual goal, time frame, budget, team, stakeholders, and so forth.
    }
    \mytext{
    
        \begin{itemize}
            \item wishes, ideas or ideals 
            \item vague goals or directions
            \item positions or functions
        \end{itemize}
    
        \medskip
    
        $\rightarrow$ projects need specifics 
    }
\end{frame}    

\begin{frame}{Projects are \textbf{not}...}
    \mycomment{
     As we establish a goal, we project (verb!) something into the future. This goal (or product) must not be confused with the process to achieve this goal. To achieve this goal, a plan is required. So the project is at best the execution of a complex plan to achive a specific goal.
    }
    \mytext{
     
        \colorbox{red!20}{\LARGE products $\neq$ PROJECTS!}
    
        \medskip
    
        % \begin{itemize}
        %     \item The anti-malaria product
        %     \item Mr. Smith and his qualities
        %     \item The presidency of Mr.\ Smith
        %     \item ...
        % \end{itemize}
    
        A project is \textbf{a planned process,} and \textbf{its result} is a product!
    
        \medskip
    
        So: how to get from the \emph{projected} result to the \emph{project?} 
    }
\end{frame}





\begin{frame}{Definition PMI}
    \mytext{
    
    \begin{quote}
        % \singlespacing
        \emph{All projects are a temporary effort to create value through a unique product, service or result. All projects have a beginning and an end. They have a team, a budget, a schedule and a set of expectations the team needs to meet. Each project is unique and differs from routine operations—the ongoing activities of an organization—because projects reach a conclusion once the goal is achieved.}
        
        \emph{The changing nature of work due to technological advances, globalization and other factors means that, increasingly, work is organized around projects with teams being brought together based on the skills needed for specific tasks.}
        \par
        \hfill ---Project Management Institute \cite{PMIWeb}
    \end{quote}
    }
\end{frame}

\begin{frame}{Definition PMI}
    \mycomment{
    The definition by PMI, often treated as canonical, can be itemised.
    }

    \mytext{
    \begin{itemize}
        \item effort to create value
        \item through a unique project, service or result
        \item beginning and end
        \item team
        \item budget
        \item schedule
        \item set of expectations
        \item closure; $\neq$ routine operations
    \end{itemize}
    }
\end{frame}
    
\begin{frame}{Definition PMI}
    \mytext{
    I like to add some points to this list:
    
    \begin{itemize}
        \item \emph{work breakdown / sub-projects}
        \item \emph{team (management)}
        \item \emph{stakeholders and contractors}
        \item \emph{governance and legal structure}
        \item \emph{technique and culture}
        \item \emph{risks / opportunities / alternatives}
        \item \emph{external effects in space and time}
        \item \emph{standards and regulations}
    \end{itemize}
    }
\end{frame}

\begin{frame}{Qualities of Projects (non-exhaustive list)}
    \mycomment{
        Projects can be classified in many ways. There is probably no perfect systematics. 
    }
\mytext{
    \begin{itemize}
        \item Large or small
        \item Complex or simple
        \item Environment: dynamic, static
        \item Clearly defined, loosely defined
        \item New or established
        \item Initialised, projected
        \item Global, local
        \item Urgent, not urgent
        \item Risky, safe
        \item High priority, lower priority
    \end{itemize}
}
\end{frame}

\begin{frame}{Types of Projects - by outcome}
    \mycomment{
    We can classify projects by their outcomes: 
    \begin{description}
        \item[research projects] we are investigating into something. This could be scientific endeavours, but also probing or investigating possibilities.
        \item[conceptualisation projects] seek to create a concept or plan,
        \item[implementation projects] implement an existing plan,
        \item[organisational implementation] seeks to set up an organisation that can achieve something;
        \item[technical implementation] has an existing organisation work to achieve a tangible goal.
    \end{description}

    
    }
    \showmytree{../shared/tikz/forest-projecttypes}
\end{frame}

\begin{frame}{Types of Projects - by outcome}
    \mycomment{
    These types of project build up upon one another. Every technical implementation requires an organisation that has to be set up first. Every implementation must follow a plan that has been conceptualised before. And to come up with a concept or plan, we first need research.

    Every technical project is also a research project etc. or must build up upon one. The opposite does not hold; a conceptualisation project might conclude that implementation is not feasible.

    It is thus a good idea to keep this different types of projects independent and introduce a review or break point between them.

    }
    
    \showmytree[1]{../shared/tikz/forest-projecttypes}
\end{frame}


\begin{frame}{Types of Projects - by owner}
    \mycomment{
    Another way to classify projects is the ownership. There are, as we will see, plenty of ways to organise project governance; and organising project governance in an intelligent way is indeed pivotal to project success. However, publicly owned and privately owned projects might follow different rules regarding finance, publication and consultation, and so forth.
    }
    \showmytikz{../shared/tikz/forest-projects-by-owner}
\end{frame}

\begin{frame}{Ownership Hierarchy}
    \mycomment{
    Ownership is rarely a single person, be it natural or legal person. Often, the \textbf{project sponsor} is who every pays for the project, and is rewarded the outcome. 
    
    The sponsor is rarely overviewing the project. More commonly, a project holder (German: Projektträger) is involved. That supervising body then has the actual project management done by the PMO, the Project Management Office. And finally, the actual work is done by other parties again, mostly contracted.

    That is not yet the whole story. There are public bodies like authorities that must be consulted or give their approval; there are consultants; there are many other parties involved at all levels. Finally, there are also the (future) users, and eyerone else affected by the project.
    }
    \showmytikz{../shared/tikz/tikz-ownership-hierarchy}

\end{frame}

\begin{frame}{The Three Constraints}
    \mycomment{
    Project Management is the art of delivering within constraints. There is more than one dimension; costs is one, but time is another, and quality a third. Understanding the \textbf{three constraints} is pivotal for success.

    The adage goes: if you want it quick, cheap, and good, choose two of these. 

    These three constraints correspond to each other. Delays are costly. On the other hand, money cannot always speed things up. And quality -- the sum of the quanlities the product ought to have -- is often non-negotiable. But then again, the scope may increase during the project (so-called \enquote{scope creep}); that will jeopadise both costs and time. 

    In the end, it is important to make stakeholders happy. Stakeholder satisfaction is ultimately the real constraint.
    }
    \showmytikz[caption={The Three Constraints (Introduction)},label=fig:threeconstraints-intro]{../shared/tikz/tikz-threeconstraints}

\end{frame}



\begin{frame}{Types of Projects - by size}
    \mycomment{
    Project come in all sorts of sizes. Indeed, only what has a certain degree of complexity can be justified to be called a project.

    Every very large project consists of large sub-projects, which consist of sub-projects, and so forth; the smallest \enquote{unit} in this sense is the individual task. 
    }
    \showmytikz{../shared/tikz/tikz-projects-by-size}
\end{frame}



\begin{frame}{Sub-Projects: Work Breakdown Structure}
    \mycomment{
    Once we break down projects into their components, we can come up with lists of sub-projects etc. all the way down to the individual tasks.

    This process of analysis is called \textbf{work breakdown.} A reasonable approach to every project is such work breakdown, that is finding a systematic answer to the question: what has to be done?

    In doing so, we find that some tasks depend on one another -- and sub-projects on other sub-projects, etc. A mere list (or a tree diagram like the WBS) does not show this satisfyingly.
    }
    \showmytikz{../shared/tikz/forest-projects-sizes-breakdown}
\end{frame}

\begin{frame}{Sub-Projects: Work Breakdown Structure}
    \mycomment{
    In project management, we typically speak of \textbf{work packages} rather than of sub-projects. A work package is typically a project (a complex task) assigned to a specific person.  
    }
    \showmytikz{../shared/tikz/forest-work-packages}
\end{frame}

\begin{frame}{Figures / Numbers: KPIs}
    \mycomment{
    Students typically underestimate costs by an order of magnitude (or two). You should research a variety of projects and their costs. Keep in mind that the price of a product (e.g. a bridge) is different to the price of a project (e.g. to build said bridge). 
    }
    \mytext{
    
    \begin{itemize}
        \item KPI Key Performance Indicator
        \item Money, time, quality, quantities...
        \item "What you can measure, you can manage"
        \item Data-driven decisions
    \end{itemize}
    
    \medskip

    When talking about anything in business,\\ provide figures / numbers!
    }
\end{frame}

\begin{frame}{Orders of Magnitude}
    \mytable{\documentclass{standalone}
\usepackage{hsbcommon}
\begin{document}
       \begin{tabular}{lrrll|l}
        Sc. & Acc. & IT & engl. &  & \textit{example} \\
        \toprule
         $10^3$ & TEUR & K & thousand &  & \phantom{laptop/computer}\\
        $10^4$ & & 10K & &  & \\
        $10^5$ & & 100K & & \emph{lakh} & \\
        \midrule
        $10^6$ & MEUR& 1M & million & \emph{Millionen} &  \\
        $10^7$ & & 10M & &  \emph{crore} & \\
        $10^8$ & & 100M & &  & \\
        \midrule
        $10^9$ &  & 1G & billion & \emph{Milliarden}  & \\
        $10^{10}$ & & 10G & &  & \\
        $10^{11}$ &  & 100G & &  & \\
        \midrule
        $10^{12}$ &  & 1T & trillion & \emph{Billionen} & \\
        \bottomrule
    \end{tabular}
\end{document}
}


\end{frame}

\begin{frame}{Types of Projects - by budget}
    \mycomment{
    Another classification for projects could be the size, that is the costs involved. Students are typically off by multiple orders of magnitude when it comes to estimating the costs of projects.

    But with this, students are not alone! As we will see, cost overruns are very common in complex projects---often also by orders of magnitude. Do yourself a favour and research a number of projects and their costs.
    }
\mytext{
    \mytable{
    
        \begin{tabular}{rl}
        \toprule
        bracket & \textit{example} \\
        \midrule
            10--50 TEUR & \\
            50--100 TEUR & \\
            100--250 TEUR & \\
            250--500 TEUR & \\
            500--1000 TEUR & \\
            1--2 MEUR & \\
            2--5 MEUR & \\
            5--10 MEUR & \\
            ... & \\
        \bottomrule
        \end{tabular}
    }

    \bigskip

    Loft bed? Drilling platform? Airport?

    Bridge? Ferry boat? Office move?

    SAP introduction? Marketing campaign?
    
}

\end{frame}

\begin{frame}{(Potential) Project Key Data}
    \mycomment{
    All projects can be described by some fundamental core data points. As project management is a business activity, you must familiarise yourself with key figures. The touch and feel, the colours and design are of less important for us here---even though for marketing and publicity, they may sometimes be more important than the hard facts!
    }
    \mytext{
    \begin{itemize}
        \item Budget: pre/post
        \begin{itemize}
            \item typically underestimated
            \item often by a factor of 2-5
            \item often by an order of magnitude
        \end{itemize}
        \item Budget variance?
        \item Schedule: duration etc.
        \item Team, Stakeholders
        \item Complexity measures
        \item Sizes, weights, measures
        \item All sorts of business key figure
    \end{itemize}
    }
\end{frame}


\begin{frame}{Project Phases, Ideally}
    \mycomment{
    Since projects have a beginning and an end---\textit{by definition}---we can distinguish a number of phases. Each of these are interesting and worth studing individually. 

    These phases are also very different to one another in both length, costs, and number of people or firms involved. 
    }
    \showmytikz{../shared/tikz/tikz-chain-projectphases-ideal}
\end{frame}

\begin{frame}{Project Phases, Realistically}
    \mycomment{
    Only rarely do projects follow the ideal flow from the first idea to the final execution. More often, there are feedback loops, and goal posts get moved. Indeed, providing for such feedback loops is important!
    }
    \showmytikz{../shared/tikz/tikz-chain-projectphases-morereal}
\end{frame}

\begin{frame}{Project Phases per Schiersmann}
    \mycomment{
    This picture shows typical phases of project phases (for change mangement projects) by \noun{Schiersmann} \cite{schiersmannthiel}. The phases follow one another, but they always relate backwards and forwards to each other, and repeat themselves.
    }
    \showmytikz{../shared/tikz/tikz-phase-oriented-project-management}
\end{frame}

\begin{frame}{PDCA-Cylce}
    \mycomment{
    Another common way of putting this is the \textbf{PDCA-Cycle}, where planning is followed by execution (or \enquote{doing}), which must always be checked; then the plan must be adjusted.

    Using the PDCA cycle is a very important tool, especially for agile projects. This already hints to the importance of circular thinking in project management.
    }
    \showmytikz[caption={PDCA Cycle (Introduction)},label=fig:pdca-cycle-intro]{../shared/tikz/tikz-PM-smartd-pdca-cycle}
\end{frame}



\begin{frame}{Project Phases (Roberts)}
    \mycomment{
    Let us now look closer at the stages of project management. 
    
    The first phase is analysis and ideation. 
    
    It is by no means clear where first ideas come from, and what the process from first ideas to first solution sketches is. Creative techniques can be applied, such as brainstorming. Undoubtedly the best ideas come from exchange between differen stakeholders.
    }
    \showmytikz[1]{../shared/tikz/tikz-chain-projectphases-details}
\end{frame}

\begin{frame}{Project Phases (Roberts)}
    \mycomment{
    Whenever first ideas have come up, it is paramount to \textbf{prototype early:} write a first sketch or proposal, draw or sketch a solution, and circulate this among potential stakeholders for further refinement. In practice, this often takes the form of \textbf{one-pagers}, short documents to present the idea to find more input from experts and potential stakeholders.
    }
    \showmytikz[1,2]{../shared/tikz/tikz-chain-projectphases-details}
\end{frame}

\begin{frame}{Project Phases (Roberts)}
    \mycomment{
    With this sketch in place, more stakeholders can be and should be and must be identified and engaged. As the idea circulates, the sketch or proposal takes more concrete shape. In particular, it becomes clear who will have which role in the project---and also who must still be approached.

    In this phase, stakeholders are typically asked to sign Letters of Intent and hence commit to the idea in an ideal, not yet contractual form. This shows further potential stakeholders the feasibility of the project.
    }
    \showmytikz[1,2,3]{../shared/tikz/tikz-chain-projectphases-details}
\end{frame}

\begin{frame}{Project Phases (Roberts)}
    \mycomment{
    With this in place, a detailed plan can be written up: the \textbf{project governance report.} This is the detailed plan sketching:
    \begin{itemize}
        \item user requirements
        \item solution design (sketch, drawing, numbers, etc.)
        \item stakeholder list and analysis
        \item work breakdown structure
        \item project plan: the time line
        \item governance: organisational chart and responsibilities
        \item identity: a name, corporate design
        \item strategy and controls
        \item risk register
        \item etc.
    \end{itemize}
    
    }
    \showmytikz[1,2,3,4]{../shared/tikz/tikz-chain-projectphases-details}
\end{frame}

\begin{frame}{Project Phases (Roberts)}
    \mycomment{
    With the Project Governance Report, contracts can be prepared and signed. The result is incorporation of some sort of \textbf{special purpose vehicle}, that is the organisation responsible for realisation of the project. It follows a kick-off event (or multiple), a ground-breaking ceremony (when physical works begin), a going public, press releases, website release and general beginning of execution of the plan.
    }
    \showmytikz[1,2,3,4,5]{../shared/tikz/tikz-chain-projectphases-details}
\end{frame}

\begin{frame}{Project Phases (Roberts)}
    \mycomment{
    With all this in place, the project is up and running. The execution, however, follows the original plan more or less strict---depending on the nature of the project. The process is documented, there are control loops, reports are written, changes documented and changes decided; all this must be embedded in internal and external communication, with accountancy, human resources management, and all other functions of running business.
    }
    \showmytikz[4,5,6]{../shared/tikz/tikz-chain-projectphases-details}
\end{frame}

\begin{frame}{Project Phases (Roberts)}
    \mycomment{
    Finally, the running project will create the projected outcome, and the project organisation must be transformed into the organisation running the project result after all goals have been reached. Accounts must be closed, invoiced be paid, responsibility handed over, and the entire operation closed and reviewed. This phase is more often than not overlooked, and bears its own particular perils.
    }
    \showmytikz[4,5,6,7]{../shared/tikz/tikz-chain-projectphases-details}
\end{frame}

\begin{frame}{Project Phases (Roberts)}
    \mycomment{
    We will roughly follow these phases within this module as we dive deeper into the issues at hand. It should, however, be remembered that project managers must permanently look ahead, adjust plans, and review prior decisions in the light of new results. Because a project is different from business as usual by the fact that we venture into something somehow unknown and never done before, no project can always be executed according to the exact plan. Flexibility is pivotal, begins in the mindset, and must be encouraged by accommodating institutions.
    }
    \showmytikz[1,2,3,4,5,6,7]{../shared/tikz/tikz-chain-projectphases-details}
\end{frame}



\begin{frame}{The Main Phases / Preliminary Gantt}
    \mytext{
    A common way to draw a project plan is a \textbf{Gantt} diagram, which shows the phases, work packages, and milestones in detail. We will speak about Gantt diagrams in more detail later.

    The following slide is an attempt to such a Gantt diagram. You see how not all Tasks have the same length. Do you think this is realistic?

    The first phase---analyse and ideate---typically takes much, much longer. To just get from the idea to a proper formulated plan (a Project Governance Report) takes more often than not as long as the rest of the project together.
    }
\end{frame}

\begin{frame}{The Main Phases / Preliminary Gantt}
    \mytable{
        \adjustbox{width=\textwidth}{
            \begin{meinganttq}{5}
                \inserttable{gantts/mains.csv}
                \showmilestones{gantts/mains.csv}
            \end{meinganttq}
        }
    }
\end{frame}


\begin{frame}{Cost Components over time}
    \mycomment{
    We have already some building blocks for project management, and understand the various phases. But so far, we are missing some hard data:
    
    How much will this project costs---in the most optimistic scenario, and in the worst case? When will these costs occur?

    Who will be working on this project, and who else will have a say, and when?

    The figure shows typical costs structures of such projects. 
    \begin{itemize}
        \item[DEVEX] Develeopment Expenditure
        \item[CAPEX] Captail Expenditure
        \item[OPEX] Operational Expenditure
    \end{itemize}
    }
        \showmytikz{../shared/tikz/tikz-costsovertime.tex}
\end{frame}




\begin{frame}{Project Governance Report}
    \mycomment{
    The Project Governance Report is the final document ready to sign off that binds the project management office. It is typically written in the later stages of the first part of project management---necessarily before the actual kick-off. 
    }

    \mytext{
    \begin{itemize}
        \item The `constitution' of a project
        \item A contract between stakeholders
        \item The actual plan, typically comprising
        \begin{itemize}
            \item user requirements ('need' or reason)
            \item detailed solution design
            \item stakeholder analysis
            \item work breakdown structure 
            \item definition of work packages and responsibilities
            \item project plan: network / BPMN plan
            \item project plan: temporal plan / Gantt diagram
            \item governance plan: corporation, controls
            \item identity and roles
            \item communication strategy 
            \item documentation, controls and reporting
            \item risk register
        \end{itemize}
    \end{itemize}
    }
\end{frame}



